\subsection{The profile coefficient and its stability}
To add the two corrections in G1, it is convenient to express them in terms of the same profile
$\nu=kP$. When composing kernels we use the uniform row and column measures, so that each sum
over a fiber is divided by the board size. Write $B=\nu-1$, $A_R=BB^\top$, $A_C=B^\top B$,
$T_X=A_X(I-A_X)^{-1}$, $U=T_RB=BT_C$, and $V=B+U$. For $X=R,C$, let
$j_X(t)=(T_X-A_X)(t,t)$, $t_X(t)=T_X(t,t)$, and $j_0=\int j_R=\int j_C$.
The cubes in double integrals are powers of the kernel values. In this notation, the functional
that enters (T) is
\begin{align*}
 \mathfrak F={}&\tfrac12j_0^2+\operatorname{tr}(T_R^2)-\tfrac32 j_0
    -\tfrac34\sum_X\int j_X^2+\langle j_R,V j_C\rangle\\
 &-\tfrac12\sum_X\langle j_X,T_Xj_X\rangle-\tfrac13\sum_X\iint T_X^3
    +2\iint BU^2+\tfrac23\iint U^3,\qquad
 \Phi=\mathfrak F-\tfrac1{12}.
\end{align*}
Let $\bar e$ be the projection of $\nu$ in $L^2(\nu)$ onto the orthogonal complement of the
additive row and column functions. Set $m_R=\int\nu^2\,dy$, $m_C=\int\nu^2\,dx$. The
deterministic part of the coefficient is
\begin{align*}
 \delta_1={}&\tfrac1{12}+\tfrac56\iint\nu^3+\tfrac12\iint\nu\bar e^2
 -\iint V\nu\bar e-\tfrac12\iint\nu^2\\
 &+\tfrac12\iint\nu^2(t_R+t_C)-\iint\nu^2V-\tfrac14\sum_X\int m_X^2 .
\end{align*}
To obtain $\delta_1$ from the normalization, we relate the stationary point $p$ to the scaling
$P$. There are potentials $\alpha,\beta$ such that
$P_{ij}=\exp(\alpha_i+\beta_j)p_{ij}(1-p_{ij})$. Multiplying the logarithm by $P_{ij}$ and
summing eliminates the gauge. Substituting in the exact expression for $\Delta$ and expanding
the entropy and the pseudodeterminant gives $\Delta=\delta_1/k+O(k^{-2})$.
The expansion uses diffusion to bound the power sums and the gap to bound the inverses.
To verify the coefficient, use $V B^\top=T_R$, $B^\top V=T_C$, $T_X A_X=T_X-A_X$, and the fact
that $B$ has zero marginals. If $g=\iint B^2$ and $\gamma_X(t)=A_X(t,t)$, the projection
above takes the form
\[
 \bar e=B+g-\rho_R-\rho_C,\quad
 \rho_R=(I+T_R)\gamma_R-V\gamma_C,\quad
 \rho_C=(I+T_C)\gamma_C-V^\top\gamma_R .
\]
Substituting this expression and $t_X=\gamma_X+j_X$ into $\delta_1$ and the static formula
gives $\mathfrak f_{\rm e}+\delta_1=\Phi$ on each grid. Hence $R=R_{\rm est}+R_{\rm det}$, with
$R_{\rm est}=\mathcal E(P)-\mathfrak f_{\rm e}/k$ and $R_{\rm det}=\Delta-\delta_1/k$.
Both are $O(k^{-2})$, proving (T). The ratio of two permanents involves the difference between
their errors. The decomposition lets us estimate this difference separately for $R_{\rm est}$ and
$R_{\rm det}$.

\paragraph{Static stability $(L^{*+})$ [D].}
For $(L^*)$ it sufficed to bound all diagrams of excess at least two. Once the board changes, we
must determine which of their contributions survive in the minor. Retain the diagrams of excess
two and three, with the correction for coincident indices $-\frac56\sum C_{aa}^3$. The same
bounds leave a tail of order $O(k^{-4})$. The gamma average and the logarithm give
$R_{\rm est}=f_2/k^2+f_3/k^3+O(k^{-4})$, with $f_2,f_3$ uniformly bounded. On passing to a
minor, the mass defect per fiber is $O(k^{-1})$. The potentials that correct it have supremum
norm $O(k^{-1})$. Removing a single square concentrates the defect on two fibers. The gauge-fixed
inverse has bounded supremum norm and off-diagonal entries of order $O(k^{-1})$. The potentials
are $O(k^{-1})$ on the two fibers and $O(k^{-2})$ elsewhere. The perturbation argument preserves
the class with bounds $2M_0,\gamma/2$. Expanding the difference of a product of propagators marks
the edge that changed. The remaining sums are bounded over a spanning tree using the absolute row
norm and the entrywise bound. Restricting the index to the affected fiber gains a factor of
$k^{-1}$ for a localized deletion. The coefficients $f_2,f_3$ change by $O(k^{-1})$ on passing
to the minor and by $O(k^{-2})$ under deletion. Accounting also for the change in the denominators
gives orders $k^{-3}$ and $k^{-4}$ for the respective increments of $R_{\rm est}$.
